\section{Conclusion}
\label{sec:conclusion}

Production LLM-agent systems are instrumented for detection, and prompt platforms
now provide real versioning, live labels, rollback, and permissions. Governance is
still uneven across the workflows, tools, MCP configurations, skills, and retrieval
corpora that determine agent behaviour. The remaining problem is therefore not an
empty release ecosystem; it is a fragmented one with different guarantees by
resource family.

\sys occupies that position. Its organizing abstraction is the \gasset: a typed
record with version history, an authority model, an audit trail, and --- where its
family supports one --- a \livetarget{} that client code resolves at call time.
Around it, \statLayersW{} layers and \statModesW{} lifecycle modes compose into a
\chainname{} carrying a flagged trace to a measured candidate, an enforced gate, an
explicit human decision, and an audited release with a recorded rollback target.

The architectural result we would most want carried forward is a capability
boundary. Adjacency in a layered stack confers capability \emph{availability}, not
\emph{inheritance}: a family obtains lifecycle behaviour and governance enforcement
only by explicitly wiring them. A meaningful base contract can still be uniform ---
identity, history, authority, and audit --- while release, rollback, evaluation, and
evidence remain declared capabilities. \sys states that guarantee surface per
family but does not yet enforce every declaration mechanically, which is a limit of
this implementation rather than an impossibility result. Two further results follow
the same discipline: naming the two dual-write
boundaries is what allows a precise, checkable audit guarantee to be stated at all,
and distinguishing containment from isolation is what keeps a security claim
falsifiable.

We are equally clear about what this paper does not establish. The quantitative
evaluation is specified and unrun; \tabref{tab:results} is empty by choice rather
than by omission, because a plausible unmeasured number is more damaging than a
visible gap. There is no user study, and the claim that governed releases are more
\emph{reviewable} is therefore argued rather than demonstrated. And the evidence is
from one implementation, which cannot show that a different factoring is impossible.

What we believe the architecture does establish is narrower and, we hope, more
durable: that the release path for non-code agent artifacts is a coherent
first-class concern, that occupying it requires an organizing abstraction the
adjacent systems do not have, and that being precise about where the resulting
guarantees stop is what separates a control plane from a dashboard with opinions.
